%----------------------------------------------------------------------------------------
% Contenido del Informe - Discusión
%----------------------------------------------------------------------------------------

\chapter{Discusión}
En conclusión, este estudio ha demostrado de manera satisfactoria la eficacia del método de Newton-Raphson para aproximar las raíces de funciones de costos no lineales con una precisión estricta de 6 cifras significativas. La perfecta coincidencia entre las implementaciones iterativas en Python y las resoluciones analíticas teóricas confirma la consistencia de los resultados obtenidos. Esto valida el uso del algoritmo numérico en problemas prácticos de optimización microeconómica aplicados al análisis de la producción agrícola, permitiendo modelar con exactitud el comportamiento de la oferta de la firma.

El desarrollo del trabajo permitió determinar con rigor analítico los umbrales operativos clave para la toma de decisiones del productor. A través de la igualación entre el Costo Marginal y las curvas de Costo Variable Medio y Costo Total Medio, se calcularon las unidades óptimas de factor variable ($x$) necesarias para alcanzar el punto de cierre y el punto de nivelación, respectivamente. Asimismo, la observación de raíces matemáticas adicionales fuera del intervalo de estudio $[1; 10]$ o con valores negativos subraya la importancia de realizar una interpretación económica cuidadosa del modelo. Estas soluciones carecen de sentido práctico en procesos productivos reales, por lo que su descarte analítico es fundamental para garantizar la aplicabilidad del análisis en la asignación eficiente de recursos.

En resumen, el método de Newton-Raphson se consolida como una herramienta robusta y de rápida convergencia para resolver ecuaciones de alto grado en ingeniería industrial y economía agraria. La metodología empleada no solo facilitó la obtención del volumen exacto de kilogramos de papa ($q$) requeridos para maximizar beneficios o cubrir costos medios, sino que también demostró la utilidad de integrar la programación en Python como soporte analítico para la verificación de criterios de convergencia y la toma de decisiones empresariales basadas en métodos numéricos.